% !TeX root = ../main.tex

\section*{Введение}

Стандартный активационный анализ данных функциональной магнитно-резонансной томографии (фМРТ) направлен на оценку изменений BOLD-сигнала (blood oxygen level dependent), связанных с условиями исследования, например, при сопоставлении интересующего и контрольного задания или при обработке двух категорий стимулов~\cite{haxby2001}. При этом физиологический интерес представляет скрытая нейрональная динамика. Известно, что BOLD-сигнал лишь косвенно отражает активность нейронов, он фиксирует локальные изменения магнитно-резонансного сигнала, вызванные гемодинамическими процессами и уровнем оксигенации крови~\cite{ogawa1990,kwong1992}. В основе формирования BOLD-ответа лежит механизм нейроваскулярного сопряжения: сдвиги в нейрональной активности последовательно меняют регионарный кровоток, объем крови и локальную концентрацию дезоксигемоглобина, что в итоге и формирует регистрируемый BOLD-отклик~\cite{friston2000,buxton2004,mandeville1999}.

В активационном анализе связь экспериментального воздействия с ожидаемым BOLD-ответом приближенно описывается функцией гемодинамического ответа (hemodynamic response function, HRF). Она описывает временной ход BOLD-ответа на кратковременное предъявление стимула: нарастание, пик и возвращение к исходному уровню с постстимульным снижением. Временной ряд предъявления стимулов каждого условия сворачивается с HRF, а результат включается в общую линейную модель (ОЛМ) в качестве регрессора~\cite{friston1998b,friston1994}. Коэффициенты $\beta$ оценивают вклад этих модельных кривых в наблюдаемый ряд, а разность $\hat\beta_A-\hat\beta_B$ образует контраст условий. 

Линейный анализ опирается на соответствие выбранной HRF линейному отклику системы и на принцип суперпозиции~\cite{monti2011}. Поскольку ответ различается между областями мозга, людьми и состояниями, фиксированный шаблон может давать ошибку даже для линейной системы~\cite{handwerker2012,lindquist2009}. Кроме того, при высокой частоте предъявления стимулов линейное приближение может оказаться недостаточно точным вследствие нелинейного характера гемодинамического ответа~\cite{polimeni2021}. Базисные функции позволяют подгонять форму HRF по данным~\cite{woolrich2004}, но не устраняют требование линейности: масштабирование стимуляции должно пропорционально масштабировать отклик, а отклики отдельных событий должны суммироваться. При сравнении условий важно также, в какой мере ошибки оценок сокращаются при вычитании коэффициентов~\cite{friston1994,wager2003}. Точная оценка контраста может сочетаться со смещенными оценками обеих амплитуд, если их смещения близки, а несократившиеся ошибки искажают оценку различия условий. Поэтому в настоящей работе точность оценок $\hat\beta_A$, $\hat\beta_B$ и их разности рассматривается раздельно.

Линейное описание может нарушаться из-за динамики сосудистой системы. Если следующее событие наступает до возвращения кровотока, объема крови и содержания дезоксигемоглобина к состоянию покоя, новый стимул действует на систему в измененном состоянии. В линейной системе перекрывающиеся отклики просто складываются, а в нелинейной измененное состояние меняет добавочный ответ на новый стимул, поэтому суперпозиция нарушается. Отклик на изолированное событие конечной амплитуды также не обязательно пропорционален малому отклику. Разреженная стимуляция ослабляет взаимодействие соседних ответов, но не гарантирует линейность одиночного отклика. При разных последовательностях предшествующих событий нелинейность способна создать ненулевую оценку контраста даже при равных амплитудах стимуляции.

Нелинейность BOLD-отклика подтверждена как экспериментально, так и в вычислительных моделях. Ответ зависит от амплитуды и частоты повторения стимулов~\cite{friston1998,robinson2006}, различается между областями мозга~\cite{birn2001}, преобразование кровотока в BOLD-сигнал нелинейно~\cite{mechelli2001}, а нейроваскулярное сопряжение само может быть нелинейным~\cite{magri2011}. Модель Balloon~\cite{buxton1998} и динамическое причинное моделирование~\cite{friston2003} описывают возможные механизмы этой нелинейности. В плотных событийных сериях различия плотности событий между условиями могут создавать искусственные различия линейных оценок и приводить к ложноположительным результатам статистического анализа~\cite{wager2005}. В настоящей работе исследуются сами ошибки оценок, а не частота таких ложных решений. При этом качество описания BOLD-ряда и точность активационных оценок не совпадают: часть нелинейного расхождения переходит в коэффициенты ОЛМ и их разность. Поэтому точность линейного анализа должна оцениваться не только по качеству описания временного ряда, но и по ошибкам $\beta$ и контраста.

Амплитуды нейрональной активности в экспериментальных данных неизвестны, поэтому прямая проверка точности этих оценок невозможна и ее проводят на синтетическом сигнале с заданным входом. В вычислительных исследованиях фМРТ~\cite{welvaert2014,ellis2020,cakan2023} такой сигнал порождает модель-генератор: BOLD получают либо интегрированием полного гемодинамического процесса~\cite{havlicek2015}, либо сверткой нейронального ряда с фиксированной HRF~\cite{sanzleon2015}. Линейный генератор по построению не воспроизводит нелинейное искажение оценок, так что для такой проверки нужен нелинейный; этот подход уже применялся для оценки смещения линейного анализа~\cite{mechelli2001,wager2005}.

Учесть нелинейность можно разными способами, но каждый требует дополнительных оценок или данных. Нелинейные методы, такие как динамическое причинное моделирование~\cite{friston2003}, подгоняют параметры по самим анализируемым данным, а эти параметры плохо идентифицируемы по BOLD-сигналу. Члены второго порядка можно оценивать и в рамках ОЛМ, но тоже по самим данным~\cite{friston1998}, а параметры нелинейности для построения регрессоров в~\cite{wager2005} оценивались в отдельном эксперименте. Другой путь не меняет исходную регрессию: нелинейный вклад рассчитывается по заранее заданной модели и однократно вычитается из уже полученных коэффициентов ОЛМ. Такая поправка не требует подгонки параметров по анализируемым данным, однако ее польза не очевидна: ошибка заданной модели, вклады высших порядков и шум могут привести к перекомпенсации. В просмотренной литературе мы не нашли оценки размера этого эффекта при разной величине контраста, несоответствии модели и шуме.

Цель работы состояла в количественной оценке ошибок коэффициентов $\beta$ и контраста $A-B$ при линейном анализе нелинейно сформированного BOLD-сигнала и размера эффекта одношаговых аналитических поправок к этим оценкам. От работ, где смещение линейного анализа исследовалось на нелинейном генераторе~\cite{mechelli2001,wager2005}, настоящее исследование отличается тремя особенностями. Поправки не добавляют членов в модель и не оценивают их по данным: квадратичное ядро получено аналитически из уравнений генератора, а рассчитанный вклад вычитается из готовых оценок. Оценка с канонической HRF сопоставляется с оценкой при собственном линейном приближении $L_1$ генератора, что отделяет вклад линейной формы от нелинейного вклада, а проекция нелинейного остатка на матрицу дизайна связывает расхождение временных рядов с ошибками оценок. Зафиксированные до проверки поправки оцениваются не по качеству описания ряда, а по ошибкам оценок относительно известного эталона, на независимых наборах условий при разной величине контраста, несоответствии модели и шуме. Постановка остается модельной: два генератора, P-DCM~\cite{havlicek2015} и Balloon--Windkessel, имеют общий линейный блок возбуждения--торможения и получают одни и те же последовательности стимулов и амплитуды, поэтому каждая ошибка известна точно. Неизвестность нейрональных амплитуд в реальных данных такая постановка не устраняет, но позволяет проверить оценки при известном ответе.